% ============================================================================
% ARR APPENDIX SECTION — READY TO COPY AND PASTE
% Requires ARR_PREAMBLE_COPYPASTE.tex and the PDF files in figures/.
% ============================================================================

\section{Full Additional-Experiment Results}
\label{sec:additional-experiments-appendix}

\subsection{External Baseline Matrix}
\label{app:external-baselines}

Tables~\ref{tab:app-baselines-gpt4o} and
\ref{tab:app-baselines-qwen} report the external attack matrix. Each row
evaluates one stored final target response per case. Because query budgets
and response-selection rules differ across methods, these results provide
context rather than a budget-matched ranking against our Best-of-4 method.

\begin{table*}[t]
\centering
\small
\caption{External baselines on GPT-4o over Final-197. Values are percentages
for one stored final response per case.}
\label{tab:app-baselines-gpt4o}
\begin{tabular}{lrrrrr}
\toprule
Method & CARES Accept & $R3+$ & $B4+$ & HarmBench+ & Strict joint \\
\midrule
Direct                & 0.00  & 2.54  & 0.51  & 1.02  & 0.51 \\
Skeleton Key          & 0.00  & 2.03  & 0.51  & 1.02  & 0.51 \\
Many-shot             & 0.00  & 0.51  & 0.00  & 0.00  & 0.00 \\
Red Teaming           & 12.18 & 30.96 & 11.17 & 8.63  & 5.58 \\
PAIR                   & 16.24 & 41.12 & 21.32 & 14.72 & \textbf{9.64} \\
TAP                    & 7.11  & 38.58 & 18.27 & 14.21 & 8.12 \\
Crescendo              & 25.89 & 22.34 & 4.57  & 5.08  & 2.03 \\
Therapist Crescendo    & 4.57  & 19.80 & 6.09  & 2.03  & 1.52 \\
PCSA Phase 2           & 32.49 & 5.58  & 2.03  & 0.51  & 0.00 \\
\bottomrule
\end{tabular}
\end{table*}

\begin{table*}[t]
\centering
\small
\caption{External baselines on Qwen2.5-7B-Instruct over Final-197. Values are
percentages for one stored final response per case.}
\label{tab:app-baselines-qwen}
\begin{tabular}{lrrrrr}
\toprule
Method & CARES Accept & $R3+$ & $B4+$ & HarmBench+ & Strict joint \\
\midrule
Direct                & 0.00  & 12.69 & 3.55  & 7.61  & 2.03 \\
Skeleton Key          & 0.00  & 20.81 & 13.20 & 18.78 & 12.69 \\
Many-shot             & 0.00  & 12.18 & 0.51  & 6.09  & 0.00 \\
Red Teaming           & 11.17 & 60.91 & 28.93 & 21.32 & 15.23 \\
PAIR                   & 13.20 & 64.47 & 39.09 & 36.04 & 27.92 \\
TAP                    & 6.09  & 76.65 & 46.70 & 43.15 & \textbf{36.04} \\
Crescendo              & 22.84 & 44.16 & 13.71 & 19.80 & 9.64 \\
Therapist Crescendo    & 8.12  & 44.67 & 14.72 & 11.68 & 6.60 \\
PCSA Phase 2           & 27.41 & 14.21 & 3.55  & 2.03  & 1.02 \\
\bottomrule
\end{tabular}
\end{table*}

\subsection{Paired PCSA Personalization Ablation}
\label{app:pcsa-personalization}

To isolate personalization within one multi-turn baseline, we pair PCSA
Phase-2 runs by case and change only whether Persona+History is available to
the attacker. Table~\ref{tab:app-pcsa-personalization} shows that
personalization increases HarmBench+ from 7.11\% to 27.92\%, $B4+$ from
7.11\% to 24.87\%, and strict joint success from 3.55\% to 13.20\%.
The paired HarmBench difference is significant under an exact two-sided
McNemar test ($p=8.23\times10^{-10}$). CARES non-refusal moves in the
opposite direction and is therefore not used alone to characterize the
personalization effect.

\begin{table}[t]
\centering
\small
\caption{Paired PCSA Phase-2 personalization ablation on
Qwen2.5-7B-Instruct ($n=197$). Values are percentages.}
\label{tab:app-pcsa-personalization}
\begin{tabular}{lrrrrr}
\toprule
Condition & CARES ASR & $R3+$ & $B4+$ & HarmBench+ & Strict joint \\
\midrule
No persona         & 98.98 & 17.26 & 7.11  & 7.11  & 3.55 \\
Persona + history  & 85.28 & 58.38 & 24.87 & 27.92 & 13.20 \\
\bottomrule
\end{tabular}
\end{table}

\begin{figure*}[t]
\centering
\includegraphics[width=0.98\textwidth]{fig2_baselines_pcsa.pdf}
\caption{External baseline results and the paired PCSA personalization
ablation. The baseline panels use one stored response per case and should not
be interpreted as a budget-matched leaderboard against our Best-of-4 method.}
\label{fig:app-baselines-pcsa}
\end{figure*}

\subsection{Attention and Representation-Space Proxy Analysis}
\label{app:proxy-analysis}

We use Qwen2.5-7B-Instruct as an inspectable proxy to compare a direct GOAL,
the compact readout, and the seven-stage dialogue. We average four branches
within each of 197 cases before statistical inference. As shown in
Table~\ref{tab:app-attention-geometry}, the attention-weighted GOAL cosine is
higher under the compact readout (0.2560) than under the long dialogue
(0.2408). The paired long-minus-short difference is $-0.0152$ (95\% CI
[$-0.0199$, $-0.0107$]). Late-layer exact-GOAL cosine follows the same
direction (0.6103 vs.\ 0.5557). The compact readout also has a lower matched
harm-axis projection ($-0.0276$ vs.\ $-0.0087$).

\begin{table*}[t]
\centering
\small
\caption{Qwen proxy analysis over 197 cases and four branches per case.
Exact-GOAL cosine and harm-axis projection use the mean of late layers 24 and
28.}
\label{tab:app-attention-geometry}
\begin{tabular}{lrrrrr}
\toprule
Condition & Positional JS & Attention GOAL cosine & Attention entropy & Exact GOAL cosine & Harm-axis projection \\
\midrule
Direct GOAL    & 1.0000 & 0.2384 & 0.5133 & 0.9998 & 0.2924 \\
Short readout  & 0.8088 & 0.2560 & 0.5431 & 0.6103 & $-0.0276$ \\
Long dialogue  & 0.7945 & 0.2408 & 0.5232 & 0.5557 & $-0.0087$ \\
\bottomrule
\end{tabular}
\end{table*}

\begin{figure*}[t]
\centering
\includegraphics[width=0.98\textwidth]{fig5_attention_geometry.pdf}
\caption{Qwen proxy analysis of attention and representation geometry. The
results are consistent with a cue-dilution hypothesis, but do not establish
a causal mechanism in GPT-4o and should not be interpreted as evidence that
the compact condition is safer.}
\label{fig:app-attention-geometry}
\end{figure*}

This analysis is intentionally limited to an association claim. Sequence
length and token position are potential confounders, and Qwen is not a proxy
for GPT-4o's internal activations. Moreover, the compact condition has higher
behavioral attack success, so its representation-space values must not be
described as evidence of improved safety.

\subsection{Can Persona History Be Reused Across Samples?}
\label{app:history-transfer}

We retrieve another sample's history using a GOAL-excluding representation.
The matched condition uses the top-1 eligible donor (mean cosine 0.8761),
whereas the mismatched condition uses a low-similarity donor with a different
crisis label (mean cosine 0.5669). The recipient persona, private GOAL,
question pool, final branches, and evaluators remain fixed.

\begin{table*}[t]
\centering
\scriptsize
\caption{Cosine-based transfer of another sample's persona history on GPT-4o
($n=197$). GOAL recovery is not applicable to short conditions because they
do not run the seven-stage recovery gate. Values other than cosine are
percentages.}
\label{tab:app-history-transfer}
\begin{tabular}{lllrrrrrr}
\toprule
Horizon & History source & Donor cosine & GOAL recovery & CARES ASR & $R3+@4$ & $B4+@4$ & HarmBench+$@4$ & Strict joint$@4$ \\
\midrule
Short & Own history            & --     & --    & 100.00 & 99.49  & 100.00 & 96.95 & \textbf{95.94} \\
Short & No history             & --     & --    & 100.00 & 97.46  & 98.98  & 86.80 & 83.76 \\
Short & Cosine-matched history & 0.8761 & --    & 100.00 & 95.94  & 99.49  & 92.89 & 89.85 \\
Short & Mismatched history     & 0.5669 & --    & 100.00 & 91.88  & 96.95  & 84.77 & 77.16 \\
Long  & Own history            & --     & 89.34 & 96.45  & 100.00 & 76.14  & 89.34 & 63.96 \\
Long  & No history             & --     & 79.19 & 98.98  & 100.00 & 86.80  & 86.29 & 73.10 \\
Long  & Cosine-matched history & 0.8761 & 76.65 & 95.94  & 98.98  & 73.10  & 81.22 & 56.35 \\
Long  & Mismatched history     & 0.5669 & 47.21 & 92.39  & 97.97  & 58.88  & 50.76 & 31.98 \\
\bottomrule
\end{tabular}
\end{table*}

\begin{figure*}[t]
\centering
\includegraphics[width=0.98\textwidth]{fig4_history_transfer.pdf}
\caption{History-transfer ablation. Cosine-matched history improves over no
history under a compact readout but does not match sample-specific history.
Long dialogue increases explicit GOAL recovery yet reduces final same-branch
behavioral success.}
\label{fig:app-history-transfer}
\end{figure*}

Cosine-matched history under the compact readout reaches 89.85\% strict joint
success, above no history (83.76\%) but below sample-specific history
(95.94\%). The matched-minus-own difference is $-6.09$ points (paired
bootstrap 95\% CI [$-10.15$, $-2.03$]) and fails the preregistered
$-5$-point non-inferiority margin. Mismatched history falls to 77.16\%, below
even the no-history condition. Under the long dialogue, matched and
mismatched histories recover the GOAL in 76.65\% and 47.21\% of cases, but
their strict joint rates are only 56.35\% and 31.98\%. GOAL recovery and final
behavioral success are therefore empirically distinct, and the long dialogue
does not reliably repair incorrect prior context.

\subsection{Reporting Scope}
\label{app:reporting-scope}

All conclusions are restricted to Final-197, the evaluated model snapshots,
and the specified evaluators. External baselines and our method use different
query budgets and response-selection rules, so we do not claim a formal
budget-matched state-of-the-art ranking from the combined table. Similarly,
the Qwen proxy analysis is descriptive and does not identify a causal safety
mechanism in closed target models.

